\documentclass[11pt,a4paper]{jsarticle}



%%%%%%%%%%%%%%%%%%%%%%%%%

% Packages

%%%%%%%%%%%%%%%%%%%%%%%%%

\usepackage{amsmath}  % 環境
\usepackage{amssymb}  % 記号
\usepackage{amsfonts} % 特殊文字

\usepackage{ascmac}

\usepackage{ulem}

\usepackage{lmodern}

\usepackage{bm}

\usepackage{mathtools, centernot, physics}

\usepackage{enumerate}

% \usepackage{graphicx}

\usepackage[dvipdfmx]{graphicx}

\usepackage{comment} % 複数行をコメントアウトする

\usepackage{url}

\usepackage{array, booktabs, multirow, makecell} % tabular関連

% \usepackage{romannum} % ローマ数字

% \usepackage{xcolor}

\usepackage[dvipdfmx, dvipsnames]{xcolor}

\usepackage{setspace}

\usepackage{float}



%%%%%%%%%%%%%%%%%%%%%%%%%

% Format

%%%%%%%%%%%%%%%%%%%%%%%%%

\setlength{\textwidth}{\fullwidth}

\setlength{\textheight}{40\baselineskip}

\addtolength{\textheight}{\topskip}

\setlength{\voffset}{-0.2in}

\setlength{\topmargin}{0pt}

\setlength{\headheight}{0pt}

\setlength{\headsep}{0pt}



%%%%%%%%%%%%%%%%%%%%%%%%%

% Definitions

%%%%%%%%%%%%%%%%%%%%%%%%%

\newcommand{\alert}[1]{\textbf{\textcolor{red}{#1}}}

\newcommand{\rainn}{\textcolor{blue}{\uline{\ \ \ \ }}}

\newcommand{\sizi}[1]{\textcolor{OliveGreen}{\lbrack #1\rbrack}}



	

\begin{document}

\section{目的}
トランジスタを用いた増幅回路として広く用いられる、電流帰還形バイアス回路によるエミッタ接地増幅回路
に関して実験を行い、その動作を理解する。

\section{実験原理}

\subsection{エミッタ接地増幅回路}

エミッタ接地増幅回路は，入力した交流信号を電圧増幅するための基本的なトランジスタ増幅回路である。トランジスタは入力信号のみでは正常な増幅動作を行えないため，あらかじめ適切な直流電圧・電流を与えて動作点を設定する必要がある。動作点が不適切であると出力波形にひずみが生じるため，直流バイアス回路は増幅回路において重要な役割を果たす。

電流帰還形バイアス回路では，エミッタ抵抗$R_E$による負帰還作用により，温度変化やトランジスタの特性変動による動作点の変化を抑制できる。ベース電圧$V_B$は抵抗$R_1$，$R_2$による分圧で与えられ，

\begin{equation}
V_B=\frac{R_1}{R_1+R_2}V_{CC}
\end{equation}

と表される。また，

\begin{equation}
V_B=V_{BE}+R_EI_C
\end{equation}

と近似できる。さらにコレクタ側では，

\begin{equation}
V_{CC}\simeq(R_C+R_E)I_C+V_{CE}
\end{equation}

が成り立ち，この式から直流負荷線が得られる。直流負荷線とトランジスタの静特性の交点が動作点となる。

\subsection{小信号等価回路}

交流信号の解析では小信号等価回路を用いる。交流解析では直流電源を交流的に接地し，結合コンデンサおよびバイパスコンデンサは十分高い周波数では短絡として扱う。トランジスタは入力抵抗$h_{ie}$と電流増幅率$h_{fe}$を用いた等価回路で表される。

このとき，電圧増幅率は

\begin{equation}
A_v=\frac{v_{\mathrm{out}}}{v_{\mathrm{in}}}
=
-\frac{R_CR_Lh_{fe}}
{(R_C+R_L)h_{ie}}
\end{equation}

となる。負号は入力信号と出力信号の位相が180$^\circ$反転することを示す。また，入力インピーダンスは

\begin{equation}
Z_{\mathrm{in}}
=
\left(
\frac{1}{R_1}
+\frac{1}{R_2}
+\frac{1}{h_{ie}}
\right)^{-1}
\end{equation}

で与えられ，出力インピーダンスは

\begin{equation}
Z_{\mathrm{out}}=R_C
\end{equation}

となる。

\section{方法}

エミッタ接地増幅回路を作製し，直流電源のみを印加してベース，コレクタおよびエミッタ各点の電圧を測定した。その後，バイパスコンデンサ$C_E$を接続した状態で交流信号を入力し，周波数を変化させながら入力電圧および出力電圧を測定して，電圧増幅率の周波数特性を求めた。

次に，$C_E$を取り外した状態で同様の測定を行い，電圧増幅率の周波数特性を測定した。

続いて，入力側に既知の抵抗を接続し，入力電圧および電流から入力インピーダンスを算出して，その周波数特性を測定した。

最後に，出力側の開放電圧と負荷抵抗接続時の出力電圧を測定し，分圧の関係から出力インピーダンスを算出して，その周波数特性を測定した。









\end{document}